\documentclass[10pt]{article}
\usepackage[margin=20mm]{geometry}
\usepackage{amsmath,amssymb,fontspec,unicode-math,xurl,hyperref}
\IfFontExistsTF{DejaVu Serif}{\setmainfont{DejaVu Serif}}{\setmainfont{Latin Modern Roman}}
\setmathfont{latinmodern-math.otf}
\setlength{\parindent}{0pt}\setlength{\parskip}{6pt}\sloppy
\begin{document}
Abstract: The public simulator does not implement the Hamiltonian claimed for the principal experiments: its kinetic coefficient is doubled and a nonzero gradient-square potential is omitted. The comparator also uses biased uniform noise in place of the specified unbiased Gaussian perturbation. These are substantive discrepancies in the current release; their impact on the historical experimental conclusions remains unresolved. Two additional algebraic defects admit local repairs. Continuous-time convergence is not disproved by those defects. No author code was executed; scientific adjudication remains pending.\par
Author: AI-assisted technical review\par
Date:\par
Title: Hamiltonian and implementation checks for gradient-based QHD\par
\# Source and scope\par
This report concerns Jiaqi Leng and Bin Shi, Quantum Optimization via Gradient-Based Hamiltonian Descent, ICML 2025, PMLR 267, OpenReview `v14jD1ankb`. The 19-page proceedings PDF was read through the main technical development and Appendices A–E, including the commutator and Lyapunov proofs. References were not independently audited. The public release https://github.com/jiaqileng/Gradient-Based-QHD, commit `1d9ed3241bd50df4d37ba6d69a8e032aef13a092`, was inspected statically: `GradBasedQHD.py`, `convex\_optim.py`, `nonconvex\_optim.py` and README. Legacy notebooks and the separate Styblinski–Tang driver were not examined. Source pages 4, 8 and 16 were actually viewed, including Figure 3. The independent SymPy check uses a one-dimensional quartic test function, not the authors’ simulator.\par
\# Problem 1: experimental implementation differs from the stated model\par
Severity: Minor. Scope: released implementation and experimental attribution; historical numerical outcomes remain unverified. Equations (13)–(14), the splitting in Section 5.1, and Appendix E.1 specify \[H=-\frac{\Delta}{2t^3}+\frac{\alpha}{2}\{-i\nabla,\nabla f\}
   +\frac{\alpha^2+\beta}{2}t^3\|\nabla f\|^2+(t^3+\gamma t^2)f.\] The convex experiment sets \(\alpha=-0.1,\beta=0,\gamma=5\) and the nonconvex experiments set \(\alpha=-0.05,\beta=0,\gamma=5\) (page 8). The release drivers call `high\_res\_qhd\_simulator` with the default `use\_new\_helper=False`, reaching `high\_res\_qhd\_helper`. Its `a1=1/t**3` is used in two Fourier factors, each with duration \(\Delta t/2\). Their combined generator is \(-\Delta/t^3\), twice the stated kinetic term. The code’s `a3=0` removes the gradient-square potential, although the required coefficient is \(0.005t^3\) in the convex experiment and \(0.00125t^3\) in the nonconvex experiments. The variable named `beta` in the release supplies the anticommutator coefficient and thus corresponds to paper \(\alpha\); a variable-name reinterpretation does not restore the missing term. The baseline `qhd\_simulator` also uses the doubled kinetic coefficient.\par
The helper approximates the anticommutator subflow with explicit Euler updates. A skew-Hermitian generator eigenvalue \(i\lambda\) gives squared norm multiplier \(|1+i\lambda\Delta t|^2=1+\lambda^2\Delta t^2\), rather than one. This is a finite-step consistency concern, not a claim about the actual error size. The code renormalizes probabilities for objective and success statistics, but `convex\_optim.py` computes gradient moments from the returned, unnormalized wavefunction. Hence those two metrics need not use the same probability normalization.\par
In addition, Appendix E.1 describes unbiased Gaussian noise in SGDM, whereas the released SGDM routine uses `np.random.rand(2)`. Each coordinate has mean \(1/2\), so the expected estimator is \(\nabla f+(1/2,1/2)\), not \(\nabla f\). This changes a principal comparator. The independent check records the coefficient ratios, missing potential coefficient and elementary norm/noise calculations. It does not execute the release or reproduce Figure 3.\par
These discrepancies affect what method and baseline the public evidence implements. They cannot establish that the paper’s actual historical figures are wrong: the authors could have used a different revision. The needed resolution is a pinned historical implementation, correct Hamiltonian coefficients and noise, normalized observables, and a time/space discretization check. Merely increasing the inner Euler count cannot fix the doubled kinetic or omitted potential. The supported assessment is Minor for the released implementation. The missing and doubled terms alter the stated algorithm, but the available evidence does not connect this commit to the historical figures or demonstrate material damage to the paper’s core empirical conclusion. Historical experimental impact remains Needs Evidence. This boundary is based on the missing provenance link, not on whether every result fails or a repair is difficult.\par
\# Problem 2: the classical Hamilton equation and ODE connection\par
Severity: Minor; local mathematical correction, no demonstrated failure of the separate quantum convergence result. Differentiating (7) with respect to \(P\) gives \[\dot X=t^{-3}P+\alpha\nabla f,\] whereas (8) prints \(P/(2t^3)\). Using the actual derivative and (9), direct differentiation cancels the Hessian–velocity terms and yields \[\ddot X+\frac3t\dot X+
 \left(1+\frac{\gamma-3\alpha}{t}\right)\nabla f
 +\beta\nabla^2 f\nabla f=0.\] Under \(\beta/\alpha=\sqrt{s}\), the last term can be rewritten as \(\sqrt{s}\nabla^2f(\dot X-P/t^3)\). The right side of the resulting high-resolution-style equation is therefore \(\sqrt{s}\nabla^2fP/t^3\), not the coefficient in (11). The independent polynomial calculation gives zero residual for the corrected equation. Both the printed first Hamilton equation and the claimed transformation require correction; one should also justify the assertion that canonical momentum decays, which is not a consequence of convergence of position alone. The paper explicitly leaves the detailed classical comparison to future work, limiting the impact of this defect.\par
\# Problem 3: a false Lyapunov identity with a valid sign repair\par
Severity: Minor. Repair verified for the differential identity on a common smooth test domain; full theorem constants are not certified. For \(\beta=0\), set \(S=\alpha t^3f\) and \(\phi=e^{iS}\Psi\). Conjugation eliminates the gradient vector potential and gives \[i\partial_t\phi=\left[-\frac{\Delta}{2t^3}
             +(t^3+\omega t^2)f\right]\phi,
 \qquad\omega=\gamma-3\alpha.\] The Lyapunov operator becomes \(O=\tfrac12\sum_j(t^{-2}p_j+2x_j)^2+(t^2+\omega t)f\). Using \(i[p_j,f]=\partial_jf\) gives the exact multiplier \[\begin{aligned}
 \partial_tO+i[H,O]
 &= (2t+\omega)f-2(t+\omega)x^\top\nabla f\\
 &=2(t+\omega)(f-x^\top\nabla f)-\omega f.\end{aligned}\] Equation (35), page 16, instead has \((2t+\omega)(f-x^\top\nabla f)-2\omega f\). For \(d=1,f=x^4/4,t=\omega=x=1\), the correct and printed values are \(-3.25\) and \(-2.75\). A symbolic differential-operator calculation on an arbitrary test function verifies the corrected expression, including all derivative terms. The discrepancy is not an OCR artifact: the page was visually checked. The exponent \(t^2\) in the denominator of (34) also differs from \(t^{1/2}\) in (33).\par
Crucially, the corrected expression remains nonpositive for convex \(f\ge0\), minimizer zero, and \(\omega\ge0\). Thus this calculation repairs the monotonicity argument instead of refuting its convergence rate. It does not independently certify the stated initial-energy bound, operator domains for merely \(C^1\) objectives, or all of Theorem 4’s parameter conditions. Those require further checking. In particular, the printed \(T_0\le1/\alpha\) condition excludes negative \(\alpha\) despite the negative values used in experiments, and Lemma 5 lists different conditions from Theorem 4; these are unresolved statement-level scope issues, not additional claimed counterexamples.\par
\# Disposition and limits\par
The core continuous-time monotonicity mechanism survives the checked local repair. The code-to-model discrepancy is substantive for empirical attribution. The spatial smoothness assumptions and complexity argument in Appendix D were read but no rigorous end-to-end quantum gate or discretization-error bound was established. Nonconvex performance, hardware advantage, and the complete experimental ranking were not reproduced. The focused independent acceptance checked all three Problems and narrowed the implementation finding to its demonstrated release-specific scope. Scientific adjudication remains pending and author adoption is unknown.\par
\# Evidence and review limits\par
The source is the frozen PMLR proceedings PDF, identified by SHA-256\par
> `845a2d3a7e6949eefce4d901a91715206c5fadecba08508e988d4778d968b5ed`\par
Physical PDF page numbers are one based and coincide with the paper’s printed page numbers in this version. The companion verification directory contains the page-indexed text extraction, the independent check program and its actual output, the coverage ledger, source and report hashes, and execution provenance. The check program was written for this review; no author program was executed. Published training runs, data collection and reported performance were not reproduced. The initial AI technical review and a fresh focused independent acceptance used GPT-6 Astra / Max selected through the platform control. Check-specific evidence and attribution are recorded in the accompanying verification. No human review or author adoption is recorded; scientific adjudication remains pending.\par
9 Jiaqi Leng; Bin Shi. Quantum Optimization via Gradient-Based Hamiltonian Descent. Proceedings of the 42nd International Conference on Machine Learning, PMLR 267 (2025). Published proceedings PDF, frozen for this review on 25 October 2026. https://proceedings.mlr.press/v267/leng25a.html.\par
\end{document}
